% Active manuscript Korean symbol index (companion to canonical v3).
\section*{현재 이론 기호 및 계산 인덱스 (한국어)}
\small
\begin{longtable}{>{\raggedright\arraybackslash}p{0.16\textwidth}>{\raggedright\arraybackslash}p{0.24\textwidth}>{\raggedright\arraybackslash}p{0.12\textwidth}>{\raggedright\arraybackslash}p{0.38\textwidth}}
\toprule
기호 & 물리적 의미 & 단위 & 정의 / 계산 방법 \\
\midrule
\endfirsthead
\toprule
기호 & 물리적 의미 & 단위 & 정의 / 계산 방법 \\
\midrule
\endhead
\bottomrule
\endfoot
$a,a_0$ & 국소 법선 층간거리와 기준 층간거리 & m & $a=a_0\lambda$. $a_0$는 별도 보정값. \\
$\lambda,\lambda_0$ & 무차원 층간거리와 기준상태 structural label & 1 & $\lambda=a/a_0$. $\lambda_0$는 $P_0$가 정의되는 기준 하중상태의 label. \\
$t,\tau,t_0$ & 물리시간, 무차원시간, 원자 기준시간 & s, 1, s & $\tau=t/t_0$, $t_0=\sqrt{m_a a_0/(EA_0)}$. reduced 초기시각은 $t=0$이며 $t_0$와 구분. \\
$m_a$ & 기준 관성질량 & kg & 미시 정상동역학 시간척도 $t_0$ 계산에 사용. \\
$E$ & 법선 강성 연결용 탄성계수 & Pa & $q=\sigma/E$, 에너지척도 $EA_ca_0$에 사용. \\
$A_0$ & 원자/기준 기계적 면적 & m$^2$ & 미시 강성 보정 기준면적. $A_c$와 동일시하지 않음. \\
$A_c$ & 국소 coherent rare event의 특성 cohesive 면적 & m$^2$ & 현재 미보정. FEM element 또는 specimen correlation 면적과 다름. \\
$m_c$ & coherent event 유효질량 & kg & 명시적 area-scaling 가정에서 $m_c=(A_c/A_0)m_a$. \\
$\sigma,q$ & 법선응력과 무차원 traction & Pa, 1 & $q(t)=\sigma(t)/E$. \\
$q_{\rm ref}$ & 기준상태 무차원 traction & 1 & 구조적 초기분포 $P_0$가 정의되는 하중수준. \\
$q_r(\lambda_0)$ & 국소 잔류 conjugate bias & 1 & $q_r(\lambda_0)=\phi'(\lambda_0)-q_{\rm ref}$. \\
$q_c$ & 최대 안정 무차원 traction & 1 & $q_c=\phi'(\lambda_c)$. \\
$m,n$ & generalized-LJ 지수 & 1 & 현재 $m=12.19$, $n=6$. \\
$\phi,\phi',\phi''$ & 무차원 상호작용에너지, traction, 접선강성 & 1 & generalized-LJ 식을 직접 계산. \\
$\lambda_c$ & 법선 불안정성 operational dividing surface & 1 & $\phi''(\lambda_c)=0$, 즉 $[(m+1)/(n+1)]^{1/(m-n)}$. 완성된 거시균열 자체를 뜻하지 않음. \\
$x_j,M$ & 미시 node 위치와 spacing 개수 & 1, 1 & finite-chain reference 변수. \\
$\lambda_i,c_i$ & 미시 label $i$의 spacing과 spacing rate & 1, $\tau^{-1}$ & $\lambda_i=x_i-x_{i-1}$, $c_i=d\lambda_i/d\tau$. \\
$E_{\rm mech}^*$ & 무차원 보존 기계에너지 & 1 & $\frac12\sum_j\dot x_j^2+\sum_i\phi(\lambda_i)$. \\
$\Gamma_0,\mu_0,\Phi^q$ & 완전 초기 미시상태, 초기상태 measure, 결정론적 flow & -- & exact microscopic push-forward에서 사용. \\
$\Lambda_i$ & exact 미시 spacing 궤적 & 1 & finite chain으로 계산되는 $\Lambda_i(\tau;\Gamma_0,q)$. \\
$F_M$ & spacing--rate empirical phase-space measure & -- & $M^{-1}\sum_i\delta(\lambda-\lambda_i)\delta(c-c_i)$. \\
$P$ & spacing 확률밀도/measure & 1 per $\lambda$ & Layer A에서는 exact projected density, reduced layer에서는 nonabsorbing structural density. \\
$u$ & 조건부 평균 spacing rate & $\tau^{-1}$ & $u=\mathrm E[c\mid\lambda,\tau]$. \\
$\Theta$ & 조건부 spacing-rate 분산 & $\tau^{-2}$ & $\Theta=\mathrm{Var}(c\mid\lambda,\tau)$. \\
$C_3$ & 조건부 3차 중심모멘트 & $\tau^{-3}$ & $C_3=\mathrm E[(c-u)^3\mid\lambda,\tau]$. \\
$\Psi$ & rate--acceleration 조건부 공분산 & $\tau^{-3}$ & $\Psi=\mathrm{Cov}(c,\ddot\lambda\mid\lambda,\tau)$. \\
$\mathcal A$ & 조건부 평균 spacing acceleration & $\tau^{-2}$ & $\mathcal A=\mathrm E[\ddot\lambda\mid\lambda,\tau]$. \\
$D_\tau$ & spacing-space 물질미분 & $\tau^{-1}$ & $D_\tau=\partial_\tau+u\partial_\lambda$. \\
$P_0(\lambda)$ & 구조적/선응력 초기 spacing 분포 & 1 per $\lambda$ & reduced model의 초기조건. thermal instantaneous width 또는 named PDF를 뜻하지 않음. \\
$\Lambda(\lambda_0,t)$ & reduced quasistatic characteristic & 1 & $\phi'(\Lambda)=\phi'(\lambda_0)+q(t)-q_{\rm ref}$의 안정 root; $\Lambda(\lambda_0,0)=\lambda_0$. \\
$\Delta\psi_c$ & $\lambda_c$까지의 무차원 유효에너지 상승량 & 1 & $[\phi(\lambda_c)-\phi'(\lambda)\lambda_c]-[\phi(\lambda)-\phi'(\lambda)\lambda]$. \\
$\Delta G_c$ & 국소 coherent event 에너지 상승량 & J & $\Delta G_c=EA_ca_0\Delta\psi_c$. \\
$T,k_B$ & 온도와 볼츠만 상수 & K, J K$^{-1}$ & transition-state factor에 사용. \\
$\nu_s$ & coherent normal coordinate의 attempt frequency & s$^{-1}$ & $\nu_s=\sqrt{\phi''}/(2\pi t_0)$. \\
$k_c$ & positive-flux transition-state escape-rate 근사 & s$^{-1}$ & $\nu_s\exp[-\Delta G_c/(k_BT)]$. high-barrier, harmonic-well, fast-equilibration 가정 및 spinodal 바로 근처가 아닌 영역에서 사용. transmission/recrossing 보정은 향후 검증 대상. \\
$P_b$ & intact survivor subdensity & 1 per $\lambda$ & $\lambda<\lambda_c$에서 quasistatic transport와 $-k_cP_b$ sink로 진화; $P_b(\lambda,0)=P_0(\lambda)$. \\
$W(\lambda_0,t)$ & characteristic 생존가중치 & 1 & $\exp[-\int_0^t k_c(\Lambda(\lambda_0,s))ds]$. \\
$S,F_{\rm ci}$ & 국소 생존확률과 누적 개시확률 & 1 & $S=\int_0^{\lambda_c}P_b d\lambda$, $F_{\rm ci}=1-S$. \\
$T_f,f$ & 피로하중 1주기 시간과 주파수 & s, Hz & $T_f=1/f$. \\
$\mathcal H_c$ & structural label의 1 cycle 누적 hazard & 1 & $\int_0^{T_f}k_c[\Lambda(\lambda_0,t),T;A_c]dt$. \\
$\mathcal H_f$ & 동일한 phase-shaped 파형을 주파수 $f$로 작동할 때의 1 cycle hazard & 1 & strict quasistatic/fast-equilibration 한계에서 $\mathcal H_f=f^{-1}\int_0^1 k_c[\Lambda_*(\theta),T]d\theta$, 따라서 $\mathcal H_f\propto1/f$. \\
$N$ & 반복 하중 cycle 수 & 1 & 동일한 loading cycle의 반복 횟수. \\
$N_{50}$ & 국소 median cycle 수 & cycle & 해당 reduced regime에서 $S_{N_{50}}=1/2$로 정의. strict fast-equilibrium 한계에서는 $N_{50}\propto f$. \\
$S_N$ & $N$ cycle 후 국소 생존확률 & 1 & mechanically stable label에 대해 $\int P_0e^{-N\mathcal H_c}d\lambda_0$. \\
$A_c/A_0$ & 특성면적 비 & 1 & 현재 민감도 분석용이며 채택값 없음. \\
$\dot D_{\rm irr}$ & 비가역 소산률 & J s$^{-1}$ & 보존 미시 baseline에서는 0; first passage와 동일시하지 않음. \\
\end{longtable}
\normalsize
